\providecommand{\Yplus}{Y^{+}}%
\providecommand{\Gquant}{G}%

\chapter{The Bulk: Three Faces of One Object}
\label{ch:the-bulk-three-faces}

The universal closed sector attached to the level-$2$
chart-specialised chiral algebra
$A_X = \Phi^{(\Sigma_{d-1}, C)}_d(\cC_X)$ on the reference curve~$C$
is the level-$3$ object identified by three constructions: the derived
chiral centre, the positive-half Hall presentation, and the Drinfeld
double of the positive half. The three meet at one object on the toric
stratum (Theorem~\ref{thm:bulk-three-faces-coincidence}). The
identification of this universal closed sector with a \emph{physical
bulk} requires the open-closed comparison arrow
\[
 \mathrm{OCA} \colon
 \cO^{\mathrm{phys}}_{\mathrm{bulk}}(\mathcal T)
 \xrightarrow{\ \sim\ }
\cZ^{\mathrm{der}}_{\mathrm{ch}}(A)
\]
of Definition~\ref{def:oca-quasi-iso} in
Chapter~\ref{ch:cy-to-chiral}; throughout this chapter, the word
\emph{bulk} refers to the closed-sector internal object
$\cZ^{\mathrm{der}}_{\mathrm{ch}}(A)$, and any physical-bulk
identification is read through OCA.

\section{Three constructions, one bulk}
\label{sec:bulk-three-constructions}

\begin{definition}[Derived chiral centre]
\label{def:bulk-zder-ch}
For $A$ an $\Eone$-chiral algebra on a smooth curve $C$, write
\[
 \cZ^{\mathrm{der}}_{\mathrm{ch}}(A)
 \;=\; \mathrm{ChirHoch}^{\bullet}(A, A)
 \;=\; \RHom_{A\text{-bimod}^{\mathrm{ch}}}(A, A),
\]
the chiral Hochschild cochain object. The bar/twisting comparison
identifies it with $\RHom(\Omega B(A), A)$ on the Koszul locus; the
identification is the comparison arrow between levels $2$ and $3$
(\S\,\ref{sec:bulk-bar-twisting}).
\end{definition}

\begin{definition}[Positive-half Hall presentation]
\label{def:bulk-yplus}
For a Calabi--Yau target $X$ with critical superpotential
$W \colon \mathcal M(X) \to \mathbb A^1$ and equivariance datum
$T_X \subset \Aut(X)$ stabilising~$W$, the \emph{positive half} is
\[
 Y^{+}(X) \;=\; H^{\bullet}_{T_X}\!\bigl(\mathcal M^{+}_{\mathrm{eff}}(X),
 \phi_{W}\bigr),
\]
the $T_X$-equivariant cohomology of the effective stable cone with
vanishing-cycle sheaf coefficients. The Hall multiplication is the
correspondence pushforward; per the four equivariance strata
(toric $T^d$, reduced $\mathbb C^{\times}\!+\!\Aut$, orbifold inertia
$I(X/G)$, lattice-polarised period), the cohomology is computed in the
ambient appropriate to the stratum.
\end{definition}

\begin{definition}[Quantum vertex group]
\label{def:bulk-gx}
The \emph{quantum vertex group} attached to the chart $(X,
\Sigma_{d-1}, C)$ is the Drinfeld double
\[
 \Gquant(X) \;=\; \mathcal D_{\hbar}\!\bigl(Y^{+}(X)\bigr),
\]
formed by adjoining a negative half $Y^{-}(X)$ and a Cartan part along a
nondegenerate Hopf pairing. The double is constructed under the
positive-half existence and the Hopf-pairing existence; both are
hypotheses on the chart, not on $X$ alone.
\end{definition}

\begin{definition}[Hall-to-chiral Morita comparison datum;
\ClaimStatusDefinitional]
\label{def:bulk-hall-chiral-morita-datum}
For a Stage-$2$ chiral algebra~$A$ on~$C$ and a positive Hall
half~$Y^+(X)$, a datum
$H_{\mathrm{Hall/ch}}(A,Y^+(X);C)$ consists of:
\begin{enumerate}[label=\textup{(\alph*)}]
\item an $E_1$-factorisation realization
\[
 Y_C^+(X):=\operatorname{Loc}_C\!\bigl(Y^+(X)\bigr)
\]
of the Hall algebra on the curve;
\item a perfect chiral
$(A,Y_C^+(X))$-bimodule~$P_X$ and a dual bimodule~$P_X^\vee$ with
evaluation equivalences
\[
 P_X\mathbin{\otimes_{Y_C^+(X)}}P_X^\vee\simeq A,
 \qquad
 P_X^\vee\mathbin{\otimes_A}P_X\simeq Y_C^+(X);
\]
\item compatibility of the induced monoidal Morita equivalence
\[
 \Xi_X:=-\mathbin{\otimes_A}P_X\colon
 \Rep^{\Eone}(A)\xrightarrow{\sim}
 \operatorname{Mod}\!\bigl(Y_C^+(X)\bigr)
\]
with the chosen nondegenerate Hopf pairing on~$Y^+(X)$ and with the
Yang--Zhao centre/double equivalence.
\end{enumerate}
Thus the datum compares module categories in their native ambients;
the associated algebra objects are related by chiral Morita
equivalence.
\end{definition}

\begin{theorem}[The bulk: conditional Morita comparison of three
level-$3$ constructions, toric $T^d$ stratum, ordinary chain ambient;
\ClaimStatusConditional]
\label{thm:bulk-three-faces-coincidence}
\textup{[Type signature: CY quadrant; toric Hall/chiral Morita
presentation; Beilinson level~$3$; hypothesis package
$H_{\mathrm{BZFN}}+H_{\mathrm{YZ}}+
H_{\mathrm{Hall/ch}}$, with $H_H$ used in clause~\textup{(v)}.]}
On the toric chart class with $X = \mathbb C^d$ ($d \in \{2, 3\}$),
$T^d$-equivariance, ordinary chain ambient, and $A = A_X^{(\Sigma_{d-1},
C)}$ the Stage-$2$ output:
\begin{enumerate}[label=\textup{(\roman*)}]
\item $\cZ^{\mathrm{der}}_{\mathrm{ch}}(A)$ exists as a chiral algebra on
$C$ (Definition~\ref{def:bulk-zder-ch}).
\item $Y^{+}(X)$ is the Schiffmann--Vasserot critical CoHA
(Schiffmann--Vasserot 2013, Publ.\ IHES 118, Theorem~8.2).
\item $\Gquant(X) = \mathcal D_{\hbar}(Y^{+}(X))$ exists as a Hopf
algebra; for $X = \mathbb C^3$ this is the affine Yangian
$Y(\widehat{\fgl}_1)$.
\item Assume the BZFN dualizability hypotheses, the Yang--Zhao
finiteness and nondegenerate-pairing hypotheses, and a datum
$H_{\mathrm{Hall/ch}}(A,Y^+(X);C)$. Then the three constructions meet
on representation categories through
\[
 \Rep^{\Etwo}\!\bigl(\cZ^{\mathrm{der}}_{\mathrm{ch}}(A)\bigr)
 \xleftarrow[\sim]{\ \mathrm{BZFN}\ }
 \cZ\!\bigl(\Rep^{\Eone}(A)\bigr)
 \xrightarrow[\sim]{\ \cZ(\Xi_X)\ }
 \cZ\!\bigl(\operatorname{Mod}(Y_C^+(X))\bigr)
 \xrightarrow[\sim]{\ \mathrm{YZ}\ }
 \Rep^{\Etwo}\!\bigl(\Gquant(X)\bigr).
\]
The first equivalence is the Ben-Zvi--Francis--Nadler centre formula
(Corollary~\ref{cor:zder-drinfeld}); the middle equivalence is induced
by the Morita kernel~$P_X$; the final equivalence is the
Drinfeld-double/centre theorem of Yang--Zhao 2018, Selecta Math.\ 24,
Theorem~B. These arrows identify the three representation categories
at Morita level.
\item If $A$ carries a family datum $H_H(A;S_A)$, then
\[
 \operatorname{Supp}\ChirHoch^\bullet(A)\subseteq S_A.
\]
This support calculation is intrinsic to the chiral chart.
\end{enumerate}
\end{theorem}

\begin{proof}
(i) is the derived-endomorphism construction in
Definition~\ref{def:bulk-zder-ch}.
(ii) is Schiffmann--Vasserot 2013, Theorem~8.2
for $\mathbb C^d$. (iii) is
the Drinfeld-double construction of a CoHA bialgebra under the stated
pairing hypotheses (Yang--Zhao 2018, Selecta Math.\ 24, Theorem~B).
For~(iv), Ben-Zvi--Francis--Nadler 2010, Adv.\ Math.\ 226, gives the
first equivalence. The invertible kernel~$P_X$ gives the induced
equivalence of Drinfeld centres. Yang--Zhao gives the final
centre/double equivalence. Clause~(v) is Volume~I Theorem~H applied to
the independent family datum $H_H(A;S_A)$.

The toric stratum and ordinary chain ambient enter the source
theorems: the equivariant cohomology in~(ii) is computed in the
$T^d$-fixed-point local model. Compact strata use the comparison
obligations of \S\,\ref{sec:bulk-compact-coha-gates}.
\end{proof}

\begin{remark}[Three constructions, three associativity classes]
\label{rem:bulk-three-associativity-classes}
The three faces carry distinct associativity classes:
\begin{itemize}[leftmargin=1.6em,itemsep=1pt]
\item $\cZ^{\mathrm{der}}_{\mathrm{ch}}(A)$ is $\Etwo$ on the centre side
(higher Deligne; Lurie 2017, HA~5.3.1.14).
\item $Y^{+}(X)$ is $\Eone$-associative cohomological-Hecke
(Schiffmann--Vasserot 2013, \S~8).
\item $\Gquant(X) = \mathcal D_{\hbar}(Y^{+}(X))$ is Hopf
(Drinfeld 1986, Dokl.~Akad.~Nauk~SSSR~283); the centre of its
representation category is $\Etwo$-braided.
\end{itemize}
The three equalities of Theorem~\ref{thm:bulk-three-faces-coincidence}
intertwine these structures by named arrows, never by elision; the
three associativity classes survive the equality.
\end{remark}

\section{The bar/twisting comparison: bar at the chart-shadow level, bulk at the next level}
\label{sec:bulk-bar-twisting}

The bar coalgebra $B(A)$ is the chart-shadow's twisting datum, living
at the same level as $A$. The bulk
$\cZ^{\mathrm{der}}_{\mathrm{ch}}(A)$ sits at the next level of the
universal arrow, with the comparison arrow given by bar--cobar
inversion composed with the chiral Hochschild identification:
\begin{equation}
 \label{eq:bulk-bar-twisting}
 \cZ^{\mathrm{der}}_{\mathrm{ch}}(A)
 \;=\; \RHom_{A\text{-bimod}^{\mathrm{ch}}}(A, A)
 \;\simeq\; \RHom\bigl(\Omega B(A), A\bigr).
\end{equation}
Equation~\eqref{eq:bulk-bar-twisting} is the level-$2$-to-level-$3$
comparison arrow under the bar--cobar theorem on the Koszul locus and
the chiral Positselski theorem in the weight-completed coderived
ambient.

\begin{remark}[Five objects: $A$, $B(A)$, $A^{i}$, $A^{!}$,
$\cZ^{\mathrm{der}}_{\mathrm{ch}}(A)$]
\label{rem:bulk-five-objects}
The five attached to an $\Eone$-chiral $A$ live at distinct loci:
$A$ at level $2$ (the chart-shadow); $B(A) = T^{c}(s^{-1}\bar A)$ at
level $2$ as the twisting coalgebra; $A^{i} = H^{\bullet}(B(A))$ the bar
cohomology coalgebra; $A^{!} = D_{\Ran}(B(A))$ the chiral Koszul dual,
the defect algebra controlling line operators ending on the boundary
where $A$ lives; and
$\cZ^{\mathrm{der}}_{\mathrm{ch}}(A)$ the bulk algebra at level $3$.
$A^{!}$ and $\cZ^{\mathrm{der}}_{\mathrm{ch}}(A)$ are distinct: at
the Kac--Moody family $A = V_{k}(\fg)$,
$A^{!} = V_{-k - 2 h^{\vee}}(\fg)$ (Feigin--Frenkel 1992, Comm.\ Math.\
Phys.\ 147) is the chiral Koszul dual, while
$\cZ^{\mathrm{der}}_{\mathrm{ch}}(A)$ is the chiral Hochschild cochain
complex. The bar represents twisting/coupling and
sits with $A$ at the chart-shadow level; the centre sits at the next
level of the universal arrow.
\end{remark}

\section{The CoHA evaluation chain}
\label{sec:bulk-coha-evaluation-chain}

For $X = \mathbb C^3$ the positive half admits the explicit evaluation
chain (Schiffmann--Vasserot 2013, Theorem~8.2; Tsymbaliuk 2017
arXiv:1703.04551, Theorem~1.1; Prochazka--Rapcak 2018,
arXiv:1812.11012):
\begin{equation}
 \label{eq:coha-evaluation-chain}
 \mathrm{CoHA}(\mathbb C^3)
 \;=\; Y^{+}(\widehat{\fgl}_1)
 \xhookrightarrow{\;\mathcal D_{\hbar}\;}
 Y(\widehat{\fgl}_1)
 \xrightarrow{\;\mathrm{ev}_{\lambda}\;}
 \End\!\bigl(\mathcal W_{1+\infty}[\lambda]\text{-vac}\bigr).
\end{equation}

\begin{theorem}[Three arrows, three associativity classes; toric
$T^3$ stratum, ordinary chain ambient]
\label{thm:coha-three-arrows}
The three arrows of \eqref{eq:coha-evaluation-chain} are distinct in
type:
\begin{enumerate}[label=\textup{(\roman*)}]
\item $\mathrm{CoHA}(\mathbb C^3) = Y^{+}(\widehat{\fgl}_1)$ is an
identification of $\Eone$-associative cohomological-Hecke algebras
(Schiffmann--Vasserot 2013, Theorem~8.2). Both sides are
\emph{positive halves}, each $\Eone$-associative bialgebras at the
Hall level.
\item $Y^{+}(\widehat{\fgl}_1) \hookrightarrow Y(\widehat{\fgl}_1)$
is the Drinfeld-double inclusion. The Hopf structure on $Y$ extends
the $\mathbb Z_2$-graded associative bialgebra structure on $Y^+$
by adjoining the negative half and the antipode, constructed on the
double (Drinfeld 1986).
\item $Y(\widehat{\fgl}_1) \xrightarrow{\mathrm{ev}_{\lambda}}
\End(\mathcal W_{1+\infty}[\lambda]\text{-vac})$ is the evaluation
homomorphism into endomorphisms of the vacuum module of the
$\mathcal W_{1+\infty}$ vertex algebra, parametrised by
$\lambda \in \mathbb C$ (Prochazka--Rapcak 2018, \S~3). The image
realises a vertex algebra $\mathcal W_{1 + \infty}$; the source is
a Hopf algebra $Y$. The arrow carries one associativity class to the
other.
\end{enumerate}
The arrows place $\mathrm{CoHA}(\mathbb C^3)$ and
$\mathcal W_{1 + \infty}$ at distinct positions in the evaluation
chain: $\mathrm{CoHA}(\mathbb C^3) = Y^+(\widehat{\fgl}_1)$ at the
Hall positive-half level; $\mathcal W_{1 + \infty}$ at the
representation-evaluation image of the doubled Hopf algebra on the
vacuum module.
\end{theorem}

\begin{proof}
(i) is Schiffmann--Vasserot 2013, Theorem~8.2; the equivariant
$T^{3}$-cohomology of the cyclic-quiver moduli space recovers the
positive half of the affine Yangian. (ii) is Drinfeld 1986; the
Drinfeld-double construction adjoins a negative half along the Hopf
pairing supplied by Schiffmann--Vasserot's Yang--Baxter axiom. (iii) is
the evaluation construction of Tsymbaliuk 2017, Theorem~1.1, sending
the spectral parameter to a complex number, with image identified by
Prochazka--Rapcak 2018 with $\End(\mathcal W_{1+\infty}[\lambda])$ on
the vacuum module. $Y^{+}$ has neither vertex-operator OPE nor a
canonical evaluation onto a vacuum; $\mathcal W_{1+\infty}$ is a
vertex algebra with state-field map. The two carry different
associativity classes; their relation factors through the doubled
Hopf algebra and the evaluation morphism.
\end{proof}

\begin{remark}[Miki triality in the evaluation chain]
\label{rem:coha-miki-triality}
The Miki $S_{3}$-triality of $\mathcal W_{1+\infty}[\lambda]$ (Miki
2007, J.\ Phys.\ A~40) is the $\mathrm{ev}_{\lambda}$-image of an
$S_{3}$-symmetry of the shuffle kernel of $Y^{+} \otimes Y^{0} \otimes
Y^{-}$ driven by $\epsilon_{i}$-permutation (Feigin--Jimbo--Miwa--Mukhin
2016, Comm.\ Math.\ Phys.~341). The triality descends through the
chain: as a Hopf-algebra triality on the doubled algebra, as a
positive-half symmetry on $Y^{+}$, and as a factorisation-algebra
symmetry on $\Ran(\mathbb C)$.
\end{remark}

\begin{remark}[The CY$_3$ specialisation $\sum \epsilon_i = 0$ supplies Miki triality at infinite rank]
\label{rem:coha-cy3-no-truncation}
On the CY$_3$ parameter surface $\sum \epsilon_i = 0$, the qq-character
recursion of $Y^{+}(\widehat{\fgl}_1)$ continues to all degrees: the
evaluation image has $\dim \mathcal W_{1 + \infty} = \infty$. The
specialisation supplies the Miki triality on the infinite-rank
algebra; the qq-character recursion forces this infinite rank across
the entire CY$_3$ locus.
\end{remark}

\section{Compact-CoHA Recognition Gates}
\label{sec:bulk-compact-coha-gates}

For $X$ a compact non-toric Calabi--Yau, the construction of
$\Gquant(X) = \mathcal D_{\hbar}(Y^{+}(X))$ from the geometric input
is controlled by a finite-and-pro recognition datum. The finite part
tests the Hall pairing, Serre ideal, coproduct, primitive centre,
associator, and parity in each height window. The pro part tests the
completion, the inverse-limit transition maps, and the Heegner divisor
comparison. The joint vanishing of these gates is the compact non-toric
existence problem for $G(X)$.

\begin{definition}[Compact-CoHA recognition gates, $K3 \times E$ chart]
\label{def:bulk-compact-coha-gates}
Let $X = K3 \times E$ with the reduced $\mathbb C^{\times} +
\Aut(X)$ equivariance datum, and let $H$ index a height filtration on
the critical Hall datum at finite truncation. The recognition gates for
a compact Hall--Drinfeld--Borcherds comparison are:
\begin{enumerate}[label=\textup{(O\arabic*)},leftmargin=2.4em,itemsep=2pt]
\item \emph{Radical descent} $\mathcal R_{H} = 0$: the radical of the
finite Hall pairing maps isometrically onto the radical of the target
Borcherds pairing.
\item \emph{Serre / PBW kernel equality} $\mathcal S_{H} = 0$: no
extra Serre relations on the source beyond those of the target;
PBW basis on the source maps to PBW basis on the target.
\item \emph{Green-adjoint coproduct} $\mathcal D^{\Delta}_{H} = 0$: the
finite Hall coproduct equals the target Drinfeld coproduct under
the Green-adjoint identification (Green 1995, Invent.\ Math.\
120).
\item \emph{Primitive-centre reduction} $\mathcal C_{H} = 0$: the
zero-charge primitive sector reduces to the centre of the target
algebra.
\item \emph{Associator cohomology} $\mathcal A_{H} = 0$: the
Hall-side associator class equals the Siegel--Borcherds-side
associator class in $\HH^{3}_{\mathrm{Hopf}}$.
\item \emph{Parity lift} $\mathcal P_{H} = 0$: the Hall parity grading
descends to the target root parity and the signed dimensions
$\dim V_{\alpha,\bar 0}-\dim V_{\alpha,\bar 1}$ equal the Borcherds
root multiplicities in the window.
\item \emph{Completion} $\mathcal Q_{H} = 0$: the height
filtration is separated, complete, and compatible with multiplication,
coproduct, antipode, and the Hopf pairing.
\item \emph{Inverse-limit exactness} $\mathcal L_{H} = 0$: the
transition maps between height windows are Mittag--Leffler, the defect
ideals commute with transition, and the derived limit term
$\varprojlim^{1}$ vanishes.
\item \emph{Heegner comparison} $\mathcal H_{H} = 0$: the root-character
pushforward sends the Hall primitive supertrace to the Heegner divisor
multiplicity and hence to the Borcherds input coefficient.
\end{enumerate}
The first six gates are finite algebraic defects. The last three are
pro/automorphic defects. Vanishing of all nine is the recognition
datum used in Definition~\ref{def:bulk-recognition-envelope}.
\end{definition}

\begin{definition}[Recognition envelope]
\label{def:bulk-recognition-envelope}
Let
\[
 \mathfrak J_{H}^{X}
 =
 \mathcal R_{H}+\mathcal S_{H}+\mathcal D^{\Delta}_{H}
 +\mathcal C_{H}+\mathcal A_{H}+\mathcal P_{H}
\]
be the finite Hall--Borcherds defect ideal in the height-$H$ Drinfeld
double $D_{H}^{X}$. The \emph{recognition envelope} for $\Gquant(X)$ at
compact chart $X$ is the universal finite quotient
\[
 \mathcal E_{H}^{X}
 \;:=\; D_{H}^{X} \big/ \mathfrak J_{H}^{X},
\]
the Drinfeld double of the finite Hall positive half quotiented by
the finite defect ideal. The original compact Hall--Drinfeld double
$D_{H}^{X}$ is recognised as the target $\Gquant(X)$ if and only if
the quotient map kills exactly the target radical and no source
relation outside the finite defect ideal, while the completed envelopes
$\widehat{\mathcal E}_{H}^{X}$ satisfy the completion,
Mittag--Leffler, and Heegner gates
$\mathcal Q_{H}=\mathcal L_{H}=\mathcal H_{H}=0$ compatibly for all
$H$.
\end{definition}

\begin{theorem}[Nine-gate compact recognition criterion]
\label{thm:bulk-nine-gate-recognition}
Fix a compact reduced $\mathbb C^\times+\Aut$ chart \(X\), a finite
height window \(H\), a finite Hall--Drinfeld source \(D_H^X\), a target
Borcherds denominator algebra \(\mathscr B_H\) with root grading, Hopf
pairing, coproduct, parity, and Heegner divisor character, and a
root-graded candidate Hopf map
\[
 \rho_H\colon D_H^X\longrightarrow \mathscr B_H .
\]
The map \(\rho_H\) recognises \(\mathscr B_H\) in height \(H\) if and
only if the six finite gates
\[
\mathcal R_H=\mathcal S_H=\mathcal D^\Delta_H=\mathcal C_H
=\mathcal A_H=\mathcal P_H=0
\]
vanish. The compatible pro-object recognises the compact quantum vertex
group \(G(X)\) only after the three pro/automorphic gates
\[
\mathcal Q_H=\mathcal L_H=\mathcal H_H=0
\]
vanish for all \(H\), with transition maps commuting with the finite
quotients.
\end{theorem}

\begin{proof}
The finite statement is an unpacking of the structure that defines an
isomorphism of Hopf-superalgebra windows. Radical descent is necessary
and sufficient for the Hopf pairing to pass to the quotient without
creating an extra null sector. Serre/PBW equality is necessary and
sufficient for the source and target presentations to have the same
relations and the same ordered monomial basis. Green-adjoint coproduct
equality is exactly compatibility of the Hall coproduct with the
Drinfeld coproduct. Primitive-centre reduction identifies the
zero-charge primitive part with the target centre. The associator class
in \(\HH^3_{\mathrm{Hopf}}\) is the obstruction to replacing a
quasi-Hopf comparison by a Hopf comparison. The parity lift is the
condition that the root multiplicities are the target superdimensions,
not only their unsigned dimensions. These six conditions are therefore
precisely the finite Hopf-superalgebra recognition conditions.

For the compact object one must pass from windows to the completed
height tower. Separated completeness makes the algebraic operations
continuous and prevents two different completed elements from having the
same finite image. Mittag--Leffler exactness kills the derived
\(\varprojlim^1\) obstruction and makes the quotient commute with the
inverse limit. The Heegner comparison identifies the primitive
supertrace in the completed Hall object with the divisor multiplicity
of the Borcherds product, hence with the level-\(4\) denominator
character. Without these three gates, the finite window comparison
does not construct \(G(X)\); with them, the inverse limit of the finite
recognitions is the compact Hall--Drinfeld object whose denominator is
the stated Borcherds product.
\end{proof}

\begin{theorem}[Five-layer answer for the global hCS--Hall question]
\label{thm:bulk-five-case-answer}
The existence of $Y^{+}(X)$ and $\Gquant(X)$ stratifies
into five layers:
\begin{enumerate}[label=\textup{(C\arabic*)},leftmargin=2.4em,itemsep=2pt]
\item $X = \mathbb C^3$: positive half constructed
(Theorem~\ref{thm:coha-three-arrows});
$\Gquant(\mathbb C^3) = Y(\widehat{\fgl}_1)$ via Drinfeld double.
\item Toric multi-chart $X$ with finite cyclic atlas: the finite
Rees hCS--Hall comparison is constructed by face-compatible cyclic
contractions and Stokes integration of the Maurer--Cartan element
over the Ran simplex; completed Rees comparison is conditional on
the Mittag--Leffler hypothesis on the Rees filtration.
\item Compact reduced $\mathbb C^{\times} + \Aut$, $X = K3 \times E$:
finite reduced compact Hall windows and radical-quotient
Hall--Drinfeld doubles constructed heightwise (Definition~\ref{def:bulk-compact-coha-gates}).
\item Recognition envelope $\mathcal E_{H}^{X}$ constructed from the
finite source by the finite defect quotient
(Definition~\ref{def:bulk-recognition-envelope}).
\item Full unquotiented Borcherds recognition of the original double
$D_{H}^{X}$: conditional on the nine gates of
Theorem~\ref{thm:bulk-nine-gate-recognition}.
\end{enumerate}
The cases stratify by ambient (ordinary / weight-completed / pro / HS-sewing)
and chart (toric / reduced / orbifold / lattice-polarised); each case has
its own obstruction set.
\end{theorem}

\begin{proof}
Each case anchors on the cited construction or named hypothesis.
A scalar character match (Igusa $\Phi_{10}$, Gritsenko $\Delta_{5}$,
or Mittag--Leffler arithmetic on the Rees filtration) does not prove
any of the nine gates. The finite gates require source proofs:
radical descent on the bilinear forms, Serre/PBW kernel equality,
Green-adjoint identification of the coproduct, primitive-centre
reduction of the zero-charge sector, associator cohomology comparison
in $\HH^{3}_{\mathrm{Hopf}}$, and parity agreement. The compact gates
require separated completion, exact inverse limit, and Heegner divisor
comparison. Each is an independent source-side condition; a target-side
scalar match has no implication for any of them.
\end{proof}

\begin{remark}[Layer-aware notation discipline]
\label{rem:bulk-layer-aware-notation}
Within this part the symbol $Y^{+}(X)$ does not pass freely through
finite Rees, completed Rees, realised critical CoHA, and
Drinfeld-double layers. Each layer uses its own decoration:
$Y^{+,\mathrm{Rees}}_{\hbar}(X)$ for the finite Rees object,
$\widehat Y^{+}_{\hbar}(X)$ for the completed Rees object,
$Y^{+,\mathrm{crit}}_{\hbar}(X)$ for the realised critical CoHA, and
$\Gquant(X) = \mathcal D_{\hbar}(Y^{+}(X))$ for the doubled object.
Each layer carries its own construction map and hypotheses; passage
between them is never automatic.
\end{remark}

\section{The K3 Yangian: the principal $d = 2$ instance of the bulk}
\label{sec:bulk-k3-yangian-principal}

At $d = 2$ the bulk on $K3$ is the K3 Yangian. The classical
Lie-algebraic envelope is $\mathfrak{so}(4, 20)$; the
Hodge-parity $\mathbb Z/2$-super-extension when imposed is
$Y_{\hbar}(\mathfrak{so}(4 \mid 20))$.

\begin{theorem}[K3 Yangian as bulk at $d = 2$, lattice-period branch,
ordinary chain ambient]
\label{thm:bulk-k3-yangian}
For $X = K3$ on the lattice-polarised period branch, with equivariant
Maulik--Okounkov presentations only on ADE/Kummer reduced-automorphism
charts:
\begin{enumerate}[label=\textup{(\roman*)}]
\item The positive half is
$Y^{+}(K3) = H^{\bullet}(\mathcal M^{+}_{\mathrm{eff}}(K3),
\phi_{W})$ with $\mathcal M^{+}_{\mathrm{eff}}(K3) = \bigsqcup_{n \ge 0}
\mathrm{Hilb}^{n}(K3)$ and $W$ trivial; the Hall product is the
Nakajima correspondence (Nakajima 1997, Duke Math.\ J.\ 91; Schiffmann
2012, Adv.\ Math.\ 230).
\item The Drinfeld double is the K3 Yangian
$Y_{\hbar}(\mathfrak{g}_{K3})$, with $24$ Heisenberg generators
indexed by the Mukai lattice $\Lambda_{\mathrm{Muk}} = U^{\oplus 4}
\oplus E_{8}(-1)^{\oplus 2}$, Mukai-signature Serre relations, and
structure function $g_{K3}(z, w)$ of bidegree $(24, 24)$
(Theorem~\ref{thm:k3-abelian-yangian-presentation}).
\item The classical limit at $\hbar \to 0$ of the abelianised
mode-Lie algebra of $Y_{\hbar}(\mathfrak g_{K3})$ is the orthogonal
Lie algebra $\mathfrak{so}(4, 20)$ preserving the Mukai pairing;
$\mathfrak{so}(4, 20)$ is symmetric on the rank-$24$ Mukai
lattice (signature $(4, 20)$), Kac type $D_{12}$.
\item The Hodge-parity $\mathbb Z/2$-super-extension is imposed on a
Mukai Hodge polarisation
$\widetilde H(K3,\mathbb R)=V_+\oplus V_-$, where $V_+$ is the
four-dimensional span of $H^0\oplus H^4$ and the real
$H^{2,0}\oplus H^{0,2}$ plane, and $V_-=V_+^\perp$ has dimension $20$.
It produces the non-Kac $Y_{\hbar}(\mathfrak{so}(4 \mid 20))$,
symmetric on both graded parts, conjectural in the (level~$3$, reduced
$\mathbb C^\times + \Aut$, ordinary chain) coordinates.
\end{enumerate}
\end{theorem}

\begin{remark}[Three form-preserving algebras, only the first two are envelopes]
\label{rem:bulk-k3-yangian-trichotomy}
Three form-preserving Lie algebras of rank $24$ admit confusion:
\begin{enumerate}[label=\textup{(\alph*)},leftmargin=1.6em,itemsep=1pt]
\item $\mathfrak{so}(4, 20)$, classical envelope, ungraded
symmetric indefinite, Kac type $D_{12}$.
\item $\mathfrak{so}(4 \mid 20)$, Hodge-parity
$\mathbb Z/2$-super-extension, symmetric on both parts, non-Kac.
\item $\mathfrak{osp}(4 \mid 20)$, Kac orthosymplectic, symmetric on
the even part and \emph{symplectic} on the odd part.
\end{enumerate}
The Mukai pairing is symmetric on both pieces of the
$V_+\oplus V_-$ Hodge polarisation. These pieces lie inside the even
Mukai cohomology; they are not the cohomological even and odd groups.
Thus $\mathfrak{osp}(4 \mid 20)$ is excluded as the K3 Yangian envelope:
it preserves the wrong form on the odd part. The Mukai envelope is
$\mathfrak{so}(4,20)$, with the Hodge-parity non-Kac refinement
$\mathfrak{so}(4\mid20)$ when the polarisation \(V_+\oplus V_-\) is part
of the datum.
\end{remark}

\begin{remark}[Nomenclature: ``K3 Yangian'' versus the K3 Hall--Drinfeld double on $K3 \times E$]
\label{rem:bulk-k3-yangian-vs-hall-drinfeld}
The label \emph{K3 Yangian} denotes the integrable quantisation of
$\mathfrak{g}_{K3}$ via its $24$ Heisenberg generators,
Mukai-signature Serre relations, and degree-$(24,24)$ structure
function. The compact $K3 \times E$ Hall--Drinfeld double
$\mathcal D_{\hbar}(\mathrm{CoHA}_{K3 \times E})$ is a separate object:
the BKM superalgebra $\mathfrak{g}_{\Delta_{5}}$ (Lorgat 2020,
arXiv:2004.09030; Gritsenko--Nikulin 1998, J.\ reine angew.\ Math.\
507) sits inside the Borcherds-recognition target on $K3 \times E$.
The K3 Yangian on $K3$ presents the self-mirror Yangian branch; the
$K3 \times E$ Hall--Drinfeld double presents the
generalised-Borcherds imaginary-root extension. The two facets meet
on rank-$24$ Mukai data; they are distinct objects.
\end{remark}

\section{The chain fusion conjecture}
\label{sec:bulk-chain-fusion}

The Stage-$2$ chiral output $A_{X}$ on the curve $C$ matches the
boundary algebra of an open factorisation dg-category. The match is
conjectural at the level of universal $X$; the toric and local models
provide constructed comparisons, and the \(K3\times E\) comparison remains
conditional on the framed \(d=3\) and compact Hall--Borcherds bridge data.

\begin{conjecture}[Chain fusion]
\label{conj:bulk-chain-fusion}
\ClaimStatusConjectured
For every CY$_{d}$-category $\cC_{X}$ and every chart datum
$(\Sigma_{d-1}, C)$, the Stage-$2$ chiral algebra
\[
 A_{X} \;=\; \Phi^{(\Sigma_{d-1}, C)}_{d}(\cC_{X})
\]
on the curve $C$ equals the boundary algebra
\[
 A_{b(X, \Sigma, C)}
\]
for a canonical boundary vacuum $b(X, \Sigma, C)$ in an open
factorisation dg-category on the tangential log curve
$(C, D_{C}, \tau_{C})$, where $D_{C}$ encodes the CY data's
special points (orbifold loci, fibration punctures, conifold
singularities) and $\tau_{C}$ is the tangential basepoint datum
of the tangential log curve.
\end{conjecture}

\begin{proposition}[Model cases for chain fusion]
\label{thm:bulk-chain-fusion-verified}
The following CY$_3$ charts provide the current model evidence for
Conjecture~\ref{conj:bulk-chain-fusion}.  In the toric and local cases
the boundary algebra is constructed by the cited local hCS/CoHA/NCCR
model.  The \(K3\times E\) entry is a conditional comparison target, not
an unconditional construction of the canonical boundary vacuum.
\begin{enumerate}[label=\textup{(\roman*)}]
\item $X = \mathbb C^3$, $\Sigma_2 = S^3 \subset \mathbb C^3$,
$C = \mathbb C \subset \mathbb C^3$. Boundary vacuum: the unique
$E_3$-augmented Stage-$1$ image at the toric origin.
$A_X = Y^{+}(\widehat{\fgl}_1)$ on $C$
(Theorem~\ref{thm:coha-three-arrows}); $A_{b}$ is the Costello--Li
6d hCS boundary algebra on $\mathbb C \subset \mathbb C^3$
(Costello--Li 2016, arXiv:1605.09930, Proposition~5.2 and \S~6).
The 5d hCS analogue at $d = 2$ identifies the Yangian VOA
$Y^{\mathrm{VOA}}(\fg)$ on a curve to all orders in $\hbar$ for
simply-laced $\fg$ (Costello--Yagi 2018, arXiv:1810.01970,
Theorem~1.1).
\item $X = $ local $\mathbb P^2 = \omega_{\mathbb P^2}$,
$\Sigma_2 = $ McKay quotient $[\mathbb C^3 / \mathbb Z_3]$,
$C = $ exceptional $\mathbb P^1$. Boundary vacuum: BKR derived
McKay equivalence $D^b\Coh(X) \simeq D^b\Coh^{\mathbb Z_3}(\mathbb C^3)$
(Bridgeland--King--Reid 2001, J.\ AMS~14, Theorem~1.2). $A_X = $
CoHA of the Beilinson 9-arrow quiver with cubic potential
(Ginzburg 2006, arXiv:math/0612139, \S~5).
\item $X = $ resolved conifold $\bar Y \to Y$,
$\Sigma_2 = $ Czech 2-chart, $C = \mathbb P^1$ exceptional curve.
Boundary vacuum: the Szendr\H{o}i 2-vertex global NCCR
(Szendr\H{o}i 2008, Geom.\ Top.~12, Theorem~2.3); two charts
$U_{\pm} \cong \mathbb C^3$ with Jordan-triple
$W_{\pm} = \tr(x [y, z])$. $A_X = D_T$-quiver CoHA on $\mathbb P^1$.
\item $X = K3 \times E$, $\Sigma_2 = $ K3 fibre over $\mathrm{pt} \in
E$, $C = E$. Candidate boundary vacuum: the framed Stage-$1$ image at
the elliptic base, after the framed \(d=3\) assignment and compact
Hall--Borcherds comparison data are supplied.  The expected target is the
vertex envelope of the BKM superalgebra $\mathfrak g_{\Delta_5}$ on $E$
(Lorgat 2020, arXiv:2004.09030, Theorem~1.2; Gritsenko--Nikulin
1998, J.\ reine angew.\ Math.\ 507, Theorem~3.1), while \(A_b\) is the
open-side boundary algebra on the elliptic base.
\end{enumerate}
These model cases anchor the required comparison maps.  They do not
construct the universal assignment \(b(X,\Sigma,C)\), and the
\(K3\times E\) case does not construct the compact Hall pairing,
Drinfeld double, or inverse-limit bridge by scalar agreement alone.
\end{proposition}

\begin{remark}[What chain fusion bridges]
\label{rem:bulk-chain-fusion-bridges}
Chain fusion identifies the Stage-$2$ output
$\Phi^{(\Sigma_{d-1}, C)}_{d}(\cC_{X})$ with the open-side primitive
$A_{b}$ in the open factorisation dg-category on
$(C, D_{C}, \tau_{C})$. The identification places the
chart-shadow at the same locus as the boundary algebra once the
canonical boundary vacuum and comparison map have been constructed.
Under that comparison, the universal closed sector
$\cZ^{\mathrm{der}}_{\mathrm{ch}}(A_X)$ would equal the closed sector
$\cZ^{\mathrm{der}}_{\mathrm{ch}}(A_b)$
of the universal-holography theorem (boundary $A$, universal closed sector
$\cZ^{\mathrm{der}}_{\mathrm{ch}}(A)$, interaction
$\mathsf{SC}^{\mathrm{ch}, \mathrm{top}}$-brace), supplying the
CY-side input to the universal-holography algebraic identification.
A physical-bulk identification on either side requires its own OCA
arrow per Definition~\ref{def:oca-quasi-iso}; the scalar
non-faithfulness lemma (Lemma~\ref{lem:scalar-non-faithfulness-vol3})
forbids any promotion from a scalar match.
\end{remark}

\begin{remark}[Open part of the chain fusion]
\label{rem:bulk-chain-fusion-open}
The conjectural part is the canonical existence of $b(X, \Sigma, C)$
for every CY$_{d}$-category $\cC_{X}$ and every chart $(\Sigma_{d-1},
C)$. The model cases above anchor on chart-specific boundary-algebra
comparisons; the \(K3\times E\) line remains conditional on the compact
comparison data. The general construction is the open frontier of canonical
boundary-vacuum construction for arbitrary $X$ and arbitrary chart
$(\Sigma_{d-1}, C)$.
\end{remark}

\section{Where the bulk lives in each chart class}
\label{sec:bulk-by-chart-class}

The four equivariance strata host four versions of
Definitions~\ref{def:bulk-yplus} and~\ref{def:bulk-gx}, each requiring
its own equivariant cohomology and its own existence obstructions.

\begin{enumerate}[label=\textup{(\arabic*)},leftmargin=2em,itemsep=2pt]
\item \emph{Toric $T^{d}$ stratum.} Chapter~\ref{ch:toric-coha}
realises the $\mathbb C^3$, local $\mathbb P^2$, resolved-conifold
positive halves explicitly via Schiffmann--Vasserot CoHA on each
toric chart; $\Gquant$ exists via Drinfeld double.
\item \emph{Reduced $\mathbb C^{\times} + \Aut$ stratum, K3 face.}
Chapters~\ref{ch:k3-yangian} and~\ref{ch:k3-quantum-toroidal} realise
$Y^{+}(K3)$ and $\Gquant(K3) = Y_{\hbar}(\mathfrak{g}_{K3})$ at the
abelian Lie level; the rational, trigonometric, and elliptic
deformation hierarchy is a Stage-$2$ specialisation of one
Stage-$1$ datum on $K3$.
\item \emph{Reduced $\mathbb C^{\times} + \Aut$ stratum, $K3 \times E$
face.} Chapters~\ref{ch:k3e-bkm}
and~\ref{ch:k3-chiral-bialgebra-platonic} realise the compact
Hall--Drinfeld--Borcherds recognition target
$\mathbf H_{\Delta_{5}}$; the recognition obstructions of
\S\,\ref{sec:bulk-compact-coha-gates} is the obstruction set.
\item \emph{Orbifold inertia} $I(X / G)$. McKay
$\mathbb C^{3} / \Gamma$ with $\Gamma \subset \mathrm{SU}(3)$ finite;
Mathieu $M_{24}$ siblings on K3.
\item \emph{Lattice-polarised period domain.} Compact CY$_3$ via
Borcherds lifts on signature-$(2, n)$ lattices; quintic; Z-manifolds;
abelian quotients.
\end{enumerate}

The four strata meet on Conjecture~\ref{conj:bulk-chain-fusion}: in
each stratum the chain fusion identifies the Stage-$2$ output with the
open-side boundary algebra, supplying the level-$3$ bulk on both
sides. The five-layer stratification of
Theorem~\ref{thm:bulk-five-case-answer} records the exact comparison
theorem available on each stratum.

\section{The bulk as the cross-volume bridge}
\label{sec:bulk-cross-volume}

The bar/twisting construction develops $E_{1}$--$E_{1}$ operadic
Koszul duality, the averaging map
$\mathrm{av} \colon \mathfrak g^{E_{1}} \to \mathfrak g^{\mathrm{mod}}$,
modular open-closed convolution, and the tangential log curve
$(X, D, \tau)$. The centre construction gives the algebraic
holographic HT sector: boundary $A$, bulk
$\cZ^{\mathrm{der}}_{\mathrm{ch}}(A)$, interaction
$\mathsf{SC}^{\mathrm{ch}, \mathrm{top}}$-brace. The CY-side
construction supplies the same level-$3$ object via
$\Phi^{(\Sigma_{d-1}, C)}_d$: the bulk
$\cZ^{\mathrm{der}}_{\mathrm{ch}}(A_X)$ on the curve $C$ has three
equivalent constructions (Theorem~\ref{thm:bulk-three-faces-coincidence}),
matches the centre construction under chain fusion, and admits the K3 Yangian as
the principal $d = 2$ instance.

The terminal scalar shadow at level $4$ (universal Borcherds-weight
identity, CHL ladder, Gritsenko--Cl\'ery $8$-form catalogue) is the
character of a level-$3$ object only after the corresponding Hall source
and comparison morphism have been constructed. On the CHL scope the
indices are \(N\in\{1,2,3,4,6\}\); on the Gritsenko--Cl\'ery scope the
indices are the triples \((t,N;k)\) of the dd-modular catalogue. The
promotion from level $4$ to level $3$ requires the
named comparison arrow under named hypotheses; the Igusa cusp form
$\Phi_{10} = D_{5}^{2}$ is not by itself the compact BPS Hilbert space
(Chapter~\ref{ch:cy-d-kappa-stratification} terminal-shadow
disclaimer; \cite{LorgatIgusaCuspForm}). The level-$4$
content is developed in its own right in the terminal-scalar part; the
bulk is the level-$3$ home from which level-$4$ scalars descend.

\section{Open frontiers attached to the bulk}
\label{sec:bulk-open-frontiers}

\begin{itemize}[leftmargin=1.6em,itemsep=2pt]
\item \emph{Chain fusion in general $d$.} The general construction of
$b(X, \Sigma, C)$ for arbitrary CY$_d$-data and arbitrary chart
$(\Sigma_{d-1}, C)$.
\item \emph{$\Gquant(X)$ for compact non-toric.} The nine recognition
gates of Definition~\ref{def:bulk-compact-coha-gates} and
Theorem~\ref{thm:bulk-nine-gate-recognition}: six finite
Hopf-superalgebra gates and three pro/automorphic gates.
\item \emph{Non-abelian K3 Yangian.} Lifting
$\widehat{\mathcal H}_{\mathrm{Muk}}^{\mathrm{ab}} \to \mathfrak{so}(4,
20)$ from the abelian-at-Lie level to a non-abelian Yangian envelope.
\item \emph{Hodge-parity super-extension $Y_{\hbar}(\mathfrak{so}(4
\mid 20))$.} Construction conditional on the $\mathbb Z/2$-graded
orthogonal Yangian datum (ACDFR 2003, arXiv:math/0304188).
\item \emph{Chain fusion respecting modularity.} Whether the open-side
trace + clutching at level $4$ matches the CY-side
$\kappa_{\mathrm{BKM}}(\Phi_N) = c_{N}(0)/2$ universal identity.
\end{itemize}
